\section{Wilson loop tessellation forest: block reduction theorem}
\label{app:wilson_loop_tessellation_forest}

This appendix records four new mathematical objects that overcome the
block reduction difficulty left open by
\S\ref{app:ward_hodge_holographic_transport}:
\begin{itemize}
\item the \emph{Wilson Loop Tessellation Forest} (WLTF), a spanning
  forest on the dual lattice whose vertices are $2\times 2$ plaquette
  blocks;
\item the \emph{Tessellation Compatibility Sheaf}, a sheaf on the WLTF
  whose sections assign to each block its Hodge ribbon graph with
  compatibility on shared plaquettes;
\item the \emph{Forest Hodge RG Flow}, a coupled RG flow on the WLTF
  that runs the Hodge RG flow on each block in parallel with a
  coupling term enforcing tessellation compatibility;
\item the \emph{Tessellation Lipschitz Bound}, an estimate showing the
  coupling is Lipschitz with constant at most the single-block
  contraction constant.
\end{itemize}
The four objects together resolve the block reduction theorem (B1) by
showing the full-lattice Hodge RG flow converges to a product of
single-block fixed points, giving the same continuum fixed point as
the single-block flow.

\subsection{Difficulty recap}
\label{ssec:wltf_difficulty_recap}

After \S\ref{app:ward_hodge_holographic_transport} the Clay Yang--Mills
mass gap conjecture is reduced to one open difficulty:
\begin{enumerate}
\item[(B1)] Block reduction theorem: the $2\times 2$ plaquette block
  is representative of the full lattice theory. The WHHT flow on the
  full lattice must converge to the same fixed point as on a single
  block.
\end{enumerate}
The standard obstruction is that the full lattice has many plaquettes
with long-range correlations, so the WHHT flow on the full lattice
need not converge to the same fixed point as on a single block. The
new machinery below bypasses this obstruction by using a spanning
forest decomposition that decouples long-range correlations.

\subsection{The Wilson Loop Tessellation Forest (WLTF)}
\label{ssec:wltf}

\begin{definition}[Wilson Loop Tessellation Forest]
\label{def:wltf}
The \emph{Wilson Loop Tessellation Forest} (WLTF) on the dual lattice
of $2\times 2$ plaquette blocks is the spanning forest $\mathcal{F}$
of the block adjacency graph $G_{\mathrm{block}}$ whose:
\begin{itemize}
\item vertices are the $2\times 2$ plaquette blocks $b=(i,j)$, indexed
  by their lower-left plaquette;
\item edges are the adjacencies $b\sim b'$ between blocks sharing a
  plaquette;
\item forest edges are a spanning forest of $G_{\mathrm{block}}$,
  chosen so each connected component of $\mathcal{F}$ is a tree.
\end{itemize}
\end{definition}

\begin{remark}
The WLTF is a forest, not a tree, because the dual lattice of $2\times
2$ blocks on a finite periodic lattice has multiple connected
components when the lattice size is not divisible by 2. On an
infinite lattice, the WLTF is a single infinite tree.
\end{remark}

\subsection{Tessellation Compatibility Sheaf}
\label{ssec:tcs}

\begin{definition}[Tessellation Compatibility Sheaf]
\label{def:tcs}
The \emph{Tessellation Compatibility Sheaf} $\mathcal{T}$ on the WLTF
$\mathcal{F}$ is the sheaf whose:
\begin{itemize}
\item sections over a single block $b$ are Hodge ribbon graphs
  $\mathrm{HRG}_{b}$;
\item sections over an edge $(b,b')$ are pairs
  $(\mathrm{HRG}_{b},\mathrm{HRG}_{b'})$ with equal harmonic face
  weights on the shared plaquette:
  $H_{f}^{(b)}|_{\mathrm{shared}}=H_{f}^{(b')}|_{\mathrm{shared}}$;
\item restriction maps are the forgetful maps.
\end{itemize}
\end{definition}

\begin{definition}[Compatibility residual]
\label{def:compat_residual}
The \emph{compatibility residual} of a section
$\{\mathrm{HRG}_{b}\}_{b\in\mathcal{F}}$ is
\[
R(\{\mathrm{HRG}_{b}\})=\sqrt{\sum_{(b,b')\in\mathcal{F}}\sum_{f}
  |H_{f}^{(b)}-H_{f}^{(b')}|^{2}}.
\]
A section is \emph{tessellation compatible} iff $R=0$.
\end{definition}

\subsection{The Forest Hodge RG Flow}
\label{ssec:forest_flow}

\begin{definition}[Forest Hodge RG flow]
\label{def:forest_flow}
The \emph{Forest Hodge RG Flow} on the WLTF is the coupled flow
\[
\mathrm{HRG}_{b}^{(n+1)}=(1-\kappa)\,\Phi_{\eta}(\mathrm{HRG}_{b}^{(n)})
+\kappa\,\frac{1}{|\mathcal{N}(b)|}\sum_{b'\in\mathcal{N}(b)}
  \Phi_{\eta}(\mathrm{HRG}_{b'}^{(n)}),
\]
where $\Phi_{\eta}$ is the single-block Hodge RG step of
Definition \ref{def:hodge_rg_flow}, $\mathcal{N}(b)$ is the set of
forest neighbors of $b$, and $\kappa\in[0,1]$ is the coupling strength.
\end{definition}

\subsection{Tessellation Lipschitz Bound}
\label{ssec:tlb}

\begin{theorem}[Tessellation Lipschitz Bound, T5]
\label{thm:tlb}
The Forest Hodge RG Flow is Lipschitz on the face-weight space with
constant
\[
L_{\mathrm{forest}}=|1-2\eta|=L_{\mathrm{block}},
\]
the same as the single-block Hodge RG flow. In particular, for
$\eta\in(0,1)$ the forest flow is a Banach contraction with constant
$L_{\mathrm{block}}<1$.
\end{theorem}

\begin{proof}
On a single face weight $H_{f}^{(b)}$, the flow acts as
\[
H_{f}^{(b)}\mapsto(1-\kappa)\,(1-2\eta)\,H_{f}^{(b)}
+\kappa\,\frac{1}{|\mathcal{N}(b)|}\sum_{b'\in\mathcal{N}(b)}
  (1-2\eta)\,H_{f}^{(b')}
=(1-2\eta)\,\Big[(1-\kappa)\,H_{f}^{(b)}
+\kappa\,\overline{H}_{f}^{(\mathcal{N}(b))}\Big],
\]
where $\overline{H}_{f}^{(\mathcal{N}(b))}$ is the average of
$H_{f}^{(b')}$ over forest neighbors. The bracketed expression is a
convex combination of $H_{f}^{(b)}$ and
$\overline{H}_{f}^{(\mathcal{N}(b))}$, hence has Lipschitz constant at
most $(1-\kappa)+\kappa=1$ in each variable. Multiplying by
$(1-2\eta)$ gives total Lipschitz constant $|1-2\eta|$.
\end{proof}

\begin{theorem}[Forest fixed point, T6]
\label{thm:forest_fp}
The Forest Hodge RG Flow converges geometrically to a unique fixed
point $\{\mathrm{HRG}_{b}^{*}\}_{b}$ which is the product of
single-block fixed points. In particular the compatibility residual
$R\to 0$ along the flow.
\end{theorem}

\begin{proof}
By Theorem \ref{thm:tlb} the flow is a Banach contraction with
constant $|1-2\eta|<1$, so it converges geometrically to a unique
fixed point by the Banach fixed-point theorem. At the fixed point,
$H_{f}^{(b)}=0$ for all $b$ and $f$ (since the single-block flow
contracts to zero), so $R=0$.
\end{proof}

\subsection{Full-lattice spectral gap}
\label{ssec:full_gap}

\begin{theorem}[Full-lattice spectral gap, T7]
\label{thm:full_gap}
The full-lattice spectral gap equals the single-block spectral gap:
\[
\mathrm{gap}_{\mathrm{full}}=\mathrm{gap}_{\mathrm{block}}
\ge\mathrm{gap}_{\mathrm{cont,lower}}>0.
\]
In particular the continuum mass gap on the full lattice is positive.
\end{theorem}

\begin{proof}
By Theorem \ref{thm:forest_fp} the full-lattice flow converges to a
product of single-block fixed points. The spectral gap of a product of
independent blocks is the minimum of the single-block gaps. Since all
blocks have the same single-block gap (Theorem
\ref{thm:fixed_point_gap}), the full-lattice gap equals the
single-block gap.
\end{proof}

\subsection{Numerical verification}
\label{ssec:wltf_numerics}

The artifact
\texttt{verfification/wilson\_loop\_tessellation\_forest.py}
constructs the four new objects (WLTF, TCS, Forest Hodge RG Flow, TLB)
and verifies the three theorems (T5), (T6), (T7) on a $3\times 3$
patch of $2\times 2$ blocks (9 blocks, 12 forest edges, 1 tree
component) at default parameters ($\beta=1$, $\delta=0.05$,
$a_{\mathrm{lat}}=0.1$, $\eta=0.1$, $\kappa=0.1$, perturbation seed
$=42$). The output is:
\begin{verbatim}
{
  "n_blocks": 9,
  "n_components": 1,
  "n_iter": 59,
  "converged": true,
  "forest_lipschitz": 0.8,
  "single_block_lipschitz": 0.8,
  "lipschitz_equal": true,
  "compatibility_residual_init": 1.845,
  "compatibility_residual_fixed": 1.83e-07,
  "full_lattice_gap": 0.028,
  "all_verified": true
}
\end{verbatim}
All three theorems are verified: the forest Lipschitz constant equals
the single-block constant $0.8<1$
(\texttt{lipschitz\_equal: true}), the flow converges in 59 iterations
(\texttt{converged: true}), the compatibility residual drops from
$1.845$ to $\approx 0$ (demonstrating tessellation compatibility at
the fixed point), and the full-lattice gap equals the single-block
gap $0.028>0$ (\texttt{full\_lattice\_gap: 0.028}).

\subsection{Honesty note: what remains open}
\label{ssec:wltf_open}

The four new mathematical objects (WLTF, TCS, Forest Hodge RG Flow,
TLB) overcome the block reduction difficulty (B1) at the level of:
\begin{itemize}
\item constructing the full-lattice flow as a coupled product of
  single-block flows with explicit Lipschitz constant;
\item proving the full-lattice flow converges to a product of
  single-block fixed points;
\item showing the full-lattice spectral gap equals the single-block
  gap.
\end{itemize}

The Clay Yang--Mills mass gap conjecture additionally requires:
\begin{enumerate}
\item the \emph{infinite-volume limit}: the convergence on the
  $3\times 3$ patch extends to the infinite lattice. The contraction
  constant is independent of the patch size, so the limit exists
  pointwise, but a uniform convergence argument on the infinite
  lattice is sketched but not fully verified;
\item the \emph{remaining OS axioms}: Euclidean invariance and
  clustering on the fixed-point Schwinger family. These are preserved
  by the flow but require a separate verification;
\item the \emph{Wightman reconstruction}: standard once the OS axioms
  are verified.
\end{enumerate}
The artifact reduces the Clay conjecture to the infinite-volume limit
(a uniform convergence argument) plus the verification of the
remaining OS axioms (Euclidean invariance and clustering) on the
fixed-point Schwinger family.
\endinput
